\section{Threats to Validity}
\label{sec:threats}

% Organised as construct / internal / external / conclusion validity, one threat
% stated once. Material enumerated rather than argued -- the coverage gaps of the
% fault-injection matrix, and the measurement incidents found during collection
% -- is in the supplementary material, pointed to from the subsection it belongs
% to. check_paper_numbers.py scans paper/supplementary.tex as well as
% main.tex and sections/*.tex, so a macro that moved with that material is still
% checked for use.

\subsection{Construct validity}

\paragraph{The barrier's durability benefit is measured on an emulated fault,
not a real power cut} \Cref{sec:eval-durability} injects the fault class
\texttt{WAITAOF} defends against --- loss of the page cache --- with a
\texttt{dm-flakey} target in \texttt{drop\_writes} mode, and the result is
unambiguous ($\FlakeyAckSurvived{}$ acknowledged records survive,
$\FlakeyUnackLost{}$ unacknowledged ones do not). What remains is that
\texttt{drop\_writes} is an \emph{emulation}. It is faithful in the dimension
that matters --- writes issued after a chosen instant never reach the platter,
while everything \texttt{fsync}-ed before it does --- and unfaithful in others:
a real power loss can tear a sector mid-write, can lose writes the device
claimed to have committed to its own volatile cache, and does not stop at a
boundary we chose. Those differences make real hardware \emph{worse} than the
emulation, so the direction of the result is safe and its precision is not: we
cannot claim the acknowledged record would survive a real power cut on a device
with a lying write cache, because \texttt{fsync} there does not mean what the
barrier assumes. That is a property of the storage, which AEP inherits from
every system that trusts \texttt{fsync}. The probe also maximises the exposure
window by phase-aligning to a just-completed fsync, so $\FlakeyUnackLost{}$ is
the loss rate for a worst-placed cut rather than a uniformly-timed one.

\paragraph{The write-loss probe tests \texttt{WAITAOF}, not AEP}
\Cref{sec:eval-durability} runs two Redis keys against a device that stops
accepting writes. It does not run AEP-full and B3 through the evaluation harness
under that fault, and so it does not measure a difference in \emph{protocol}
outcomes --- duplicates, lost effects, declared ambiguity --- attributable to
write loss. What it establishes is the premise the durability argument rests on:
that an acknowledged record survives an event that destroys an unacknowledged
one, which under a process kill is false and under write loss is true.

\paragraph{We ran the extension, and the fault is deliverable while the
durability signal above it is not trustworthy} The harness's Redis now runs with
its append-only file on the flakey device, and a full session --- both arms,
thirty runs each, ten executions per run, \texttt{drop\_writes} armed at the
intent checkpoint --- was collected without a delivery failure, so
\textbf{the fault is deliverable} for a whole session on this host. The second
half is answered in a way we did not predict. \textbf{\texttt{WAITAOF} returned
success for every durability acknowledgement requested after the device had
stopped accepting writes.} \texttt{dm-flakey}'s \texttt{drop\_writes} discards
each write bio and reports success, so the kernel sees no error,
\texttt{fsync(2)} succeeds, and Redis has no way to learn otherwise. The barrier
withholds dispatch on a \emph{failed or absent} durability acknowledgement; it
was handed a successful one and dispatched, behaving exactly as specified on an
input that was false. Both arms finished at ceiling.

We state the scope narrowly. \textbf{It is not evidence that the barrier fails
to withhold under write loss}: the condition the experiment exists to create ---
the barrier meeting an acknowledgement that fails or never arrives --- was never
reached, so no run exercised withholding. \textbf{Neither is it evidence that
the barrier works.} What it establishes is that the durability guarantee rests
on an unstated premise --- that a \texttt{WAITAOF} success means the record
reached the device --- and that under storage which discards writes
\emph{silently} the premise is false and the protocol cannot detect it. Storage
that returned \texttt{EIO} would let the acknowledgement fail and put the
barrier back on the input its guarantee is defined over; this result therefore
bounds the guarantee to storage that reports its own failures, and says nothing
about storage that does not. It is one host and one Redis version (7.2.5,
\texttt{appendfsync everysec}), and we claim it for no other.

\paragraph{The write-loss probe runs several layers above hardware, and does not
need to be closer} The \texttt{dm-flakey} target sits on a loop device backed by
a file on the WSL2 root filesystem, itself a virtual disk on the host. The
layering is irrelevant to the claim, because \texttt{dm-flakey} \emph{is} the
boundary under test: the experiment asks whether a \texttt{WAITAOF}
acknowledgement corresponds to bytes that have crossed a chosen point in the
storage path before that point stops accepting writes, and everything below the
target is on the far side of it in both arms equally. What the layering
\emph{would} threaten is a claim about absolute fsync latency on real hardware,
which we do not make from this probe --- the barrier's cost comes from
\cref{tab:deployment} --- or a claim that a real power cut cannot do worse,
which we explicitly do not make above.

\paragraph{No process-level fault can exercise the barrier's durability at all}
Stated separately because a reader is most likely to have assumed the other way.
\texttt{appendfsync everysec} defers the \texttt{fsync(2)} and not the
\texttt{write(2)}, so a \texttt{SIGKILL} of Redis loses nothing that
\texttt{appendonly yes} was not already keeping ($\ProcessKillUnackLost{}$, and
$n{=}\KillCanaryN{}$ canaries in the ablation cells). \textbf{Under a process
fault the barrier's contribution is prevention (\cref{sec:eval-prevention},
$p=\UnwantedP{}$, replicated at
[\ReplicationPreventedLow{},\,\ReplicationPreventedHigh{}] effects prevented)
and not durability}, and any reading of the barrier as protecting the record
from a crashed Redis process is wrong.

\paragraph{The barrier is evaluated under crash faults only} Every prevention
measurement here delivers an F3 crash-stop of the state store, a class in which
the arms are separated by a \emph{timing} difference and nothing else: AEP-full
has one extra Redis-dependent step, and the fault must land in a window B3 has
already left. That window can be narrowed and measured (\cref{sec:eval-prevention})
but not closed, because a fault that guarantees \texttt{WAITAOF} cannot return
also blocks the two Redis calls B3 makes after the same checkpoint. One
mechanism would avoid this --- disabling the fsync at the checkpoint while
leaving every other Redis call working, which separates the arms by construction
--- and \textbf{we designed that arm and did not run it}. It is a configuration
change that makes a guarantee unavailable rather than a fault, it would
instantiate a failure class our model does not contain, and the controlled fault
had already reduced the between-session spread to \ArmASpreadMax{} count per
class. The scope this fixes is explicit: the prevention guarantee is evaluated
where durability is \emph{raced for} and lost to a crash, not where it is simply
\emph{unavailable}. A reader who needs the latter should treat it as untested.

\paragraph{The mock provider's realism} The endpoint is a purpose-built service
with a configurable delay, timeout and error distribution, and a
transactionally written ground-truth ledger. It is realistic in the dimension
the experiment turns on --- its reconciliation capability, the independent
variable of \cref{tab:outcomes} --- and unrealistic in most others: no
authentication, no rate limits, no partial writes, no schema evolution, one
mutation shape. A real legacy endpoint may also \emph{misdeclare} its
capability, which the protocol cannot detect (\cref{sec:model-endpoint}); we
validate the capability's shape, never its truthfulness.

\paragraph{The dispatch guard assumes trusted same-process code} The guard
prevents accidental or ordinary control-flow bypass through the supported
dispatch API, and its registry detects copying, reuse and scope mismatch. It is
not a cryptographic capability: Python's underscore naming,
object-construction checks and closure placement are not security boundaries,
and code running in the same interpreter can reach module internals or bypass
the API altogether. The implementation therefore enforces a protocol invariant
only under a trusted-code assumption; a hostile-extension threat model would
require a separately enforced process or service boundary.

\paragraph{B4 is not Temporal} B4 shares exactly one mechanism with a production
durable-execution engine: durable history, replay, and at-least-once activity
re-execution. It does not model heartbeating, backoff timing, task queues,
versioning, or determinism checking. Its re-execution policy is the vendor's
documented default and is quoted as such~\cite{temporal-retry}, but the decision
is our code. No claim about Temporal's throughput, latency, maturity or
correctness as a product is made or supported anywhere in this paper. B4 and B4b
\emph{bracket} the configuration space rather than survey it: a deployment could
set Maximum Attempts to 3, add a heartbeat, or make the activity idempotent ---
the first two land between our two points on the same duplicate/loss trade, and
the third leaves the problem this paper is about.

\paragraph{The engine itself was run, and it separates what B4 licenses from
what it does not} We ran the production engine against the same non-cooperative
endpoint, in the two configurations B4 and B4b instantiate, at the one crash
point both the engine and our model can be cut at. With unlimited attempts it
duplicated at \BfiveDup{} over \BfiveExec{} executions and lost effects at
\BfiveLost{}; with Maximum Attempts set to one it lost at \BfivebLost{} over
\BfivebExec{} executions and duplicated at \BfivebDup{}. Neither configuration
declared ambiguity in any of \BfiveRuns{} non-void runs (\BfiveAmb{} and
\BfivebAmb{} declarations) --- the third corner of \cref{sec:trilemma} going
unreached by the engine itself and not merely by our model of it.

\emph{The direction transfers; the magnitudes do not.} Each configuration lands
in the corner \cref{tab:related} assigns it, with exactly zero events of the
other kind. But at that same crash point B4 duplicates at \BfourAtBarrier{} over
\BfourAtBarrierExec{} executions and B4b loses at \BfourbAtBarrier{} over
\BfourbAtBarrierExec{}, roughly six times the engine's rates. \textbf{B4's rates
are our model's, not the product's}, and no rate in this paper should be read as
an engine's.

Four things bound this and none is removed by having run the engine. The
comparison is \emph{incomplete}: the engine has no position answering to
\textsc{after-intent-before-barrier}, because its pre-dispatch record is written
by the server inside one RPC, so there is no worker-side instant between the
write and its acknowledgement --- a hole exactly where B4 and B4b have data at
all three capability classes. The rates are conditional on a
\BfiveStartToClose{}\,ms Start-To-Close timeout that B4 has no analogue for; a
different timeout is a different cell, and no timing number here is any engine's
recovery latency. The interval arithmetic is asymmetric --- B4's intervals rest
on three runs against the engine's thirty, and three clusters put resampled
means on multiples of one thirtieth --- so the separation should be read from
the point estimates, which differ by roughly $0.8$, and not from the intervals
failing to overlap. And the crash window at the dispatch point is sized by the
server and by loopback rather than by the protocol, so its width is not
like-for-like with B4's.

\paragraph{Declared ambiguity is not evaluated as an operational outcome} The
paper measures how often AEP declares ambiguity and shows that it never
substitutes a silent failure for it. It does \emph{not} show that declaring is
better in practice; that would require an operator study, and we have run none.
The argument that a declared incident beats an undiscovered one is a normative
claim about failure handling, not a result of this evaluation. Relatedly, the
artifact has no escalation \emph{mechanism} (\cref{tab:nonclaims}): reaching the
terminal state pauses the execution and alerts nobody.

\paragraph{Test counts are not protocol assurance} The suite's size is evidence
about the artifact's hygiene, not about the protocol's behaviour under faults.
The evidence for the latter is the fault-injection matrix.

\subsection{Internal validity}

\paragraph{The decomposition narrows our own novelty claim}
\Cref{sec:eval-detection} attributes the detection guarantee to the pre-dispatch
record and the transition table's missing re-entry edge, not to the barrier
machinery of contribution C2. The reading follows: our most novel mechanism
serves the claim with the \emph{narrowest} evidence --- one crash point, one
host, an effect size whose host-dependence we have measured rather than
suspected --- while the claim with the strongest evidence rests on write-ahead
logging plus a state-machine property that is a few lines of Lua. Narrowest is
not weakest: within that scope the effect replicated across
\ReplicationSessions{} further sessions at
[\ReplicationPreventedLow{},\,\ReplicationPreventedHigh{}] effects prevented,
the only session-clustered interval in this paper that excludes zero. The
composition being novel was never the argument (\cref{sec:related}); what is new
is the semantics it delivers, and the ablation says which half delivers which
--- including that the expensive half is optional.

The cell is no longer one collection under one injector.
\Cref{sec:eval-prevention} leads with a \emph{controlled} fault --- the
container frozen before it is killed, landing in \ControlledLanding{}\,ms
against \UncontrolledLanding{}\,ms --- over \ArmASessions{} sessions and all
\ArmAClasses{} capability classes. AEP-full's unwanted-applied rate is
\ArmAAepRate{} against B3's \ArmABthreeRate{}, with a between-session spread of
\ArmASpreadMax{} count per class: a measured quantity with a stated spread
rather than one draw from a distribution we could only name. \textbf{It is
bounded by measurement rather than by prediction, and that cost us a
pre-registration.} We registered [\ArmABoundLow{},\,\ArmABoundHigh{}] before
collecting, derived from the injector's landing latency as a fraction of the
\texttt{WAITAOF} window, and the observed \ArmAAepRate{} fell \emph{below} it.
Both of the model's premises fail on measurement: \ModelAcksPastLanding{} of
\ModelAckN{} uncontrolled runs acknowledged after that landing had supposedly
passed and the same injector call runs \ModelInSituRatio{}$\times$ the bench
figure in situ, and the wait is not uniform --- \ModelEmpiricalWindow{} of waits
fall under \ControlledLanding{}\,ms where uniformity assigns
\ModelUniformWindow{}. Take the measured rate and its spread, not the bound.

\textbf{The residual is real and irreducible under this fault class.} AEP-full
still applied \ArmAAepApplied{} effects in \ArmAAepN{} executions, and no
crash-class fault can drive that to zero, for the reason given under construct
validity. The uncontrolled cell agrees in direction and its magnitude still
moves: \ReplicationSessions{} collections ranged
\ReplicationAepMin{}--\ReplicationAepMax{} applied effects out of
\AepKillRuns{}. Attributing that movement to kill latency did not succeed at
this design's precision --- \KillLatencyPerSession{}\,ms per session, mean
\KillLatencyMean{}\,ms, 95\% interval
[\KillLatencyLow{},\,\KillLatencyHigh{}], a half-width
\KillLatencyPrecisionRatio{} times that mean, containing zero because the
sessions disagree. \UnwantedPrevented{} is one draw from that distribution;
\ArmAAepRate{} is not.

\paragraph{The detection finding has no referent outside this artifact} The
decomposition above is established by an ablation --- AEP-full against itself
with the barrier removed --- and an ablation is internal by construction. That
is the right instrument for \emph{attributing} an effect to a mechanism, but it
cannot corroborate the resulting claim from outside. Nothing external to this
artifact establishes that a pre-dispatch record without an fsync barrier
suffices for detection: the proposition is supported by two systems we wrote,
measured by a harness we wrote, against a provider we wrote. The oracle's
independence (below) addresses whether we counted correctly, not whether the
finding generalises beyond our implementation of the pattern. That is the
structural cost of the instrument, and it applies to the paper's most
load-bearing empirical claim.

\paragraph{The matrix found defects in our own measurement, including one that
would have inflated counts} Three are on the record. A resume defect let a
re-attempted run merge stale event shards into its own log, double-counting
executions. A re-executing supervisor abandoned the lease instead of waiting for
it, and execution-state seeding was not idempotent --- both making a
\emph{baseline} look better than it is, and so biasing against our own
hypothesis. And running the project's own destructive integration suite beside a
live collection replaced the matrix's Redis underneath it, so every B4/B4b run
from that window was discarded and re-collected. All three are fixed, the
affected runs were re-collected, and none of the discarded runs contributes a
number here; the runner now refuses to start beside a host-level fault injector
and records the server's \texttt{run\_id} with every run. The supplementary
material gives each incident in full. We record them because the general point
--- that a fault-injection harness is a stronger verification tool than a unit
suite --- cuts both ways, and there may be defects it has not found.

\paragraph{Ambiguity is counted from the oracle, not from self-report} A system
that declared ambiguity liberally could otherwise score well by saying ``I do
not know'' about everything. Duplicates and lost effects are counted by
comparing the provider's independently written ledger against the system's
records; the analysis is prevented by a parsed source gate from reading live
protocol state.

\paragraph{The oracle is independent, not infallible} The provider ledger is
written outside the systems under test, and every completed run is reconciled
against it, but the provider and ledger share one implementation and one host.
One voided run demonstrates that this oracle can fail: its event log and ledger
disagreed while sibling and repeated runs agreed. The current repository retains
the aggregate record of that exclusion, and the voided raw directory must be
included in the external archive before submission. A second provider
implementation or an external observer would reduce this common-mode doubt;
neither is present.

\paragraph{One author implemented every system under comparison} B0--B2 are
small enough to inspect, each carries a machine-readable descriptor that tests
check against its implementation, and B4's re-execution policy is defended
against the vendor's own documentation in the artifact. That is mitigation, not
elimination: a baseline written by the people proposing the alternative is a
structural conflict, and the artifact is released partly so the baselines can be
re-implemented by someone else.

\subsection{External validity}

\paragraph{Single node, single trust domain} One Redis instance, no replica, no
consensus. \Cref{tab:nonclaims} lists what this forecloses. One residual
interacts with the results: an AOF rewind can un-fence a lease, since lock keys
are deliberately not barriered, and the fix is HA rather than anything local.

\paragraph{Concurrency coverage is narrow} The harness exercises crash and
takeover with two workers, so stale-owner and successor paths are covered, but
it does not sustain live contention among many workers for the same execution or
endpoint resource. The evaluation therefore supports fencing and recovery under
the measured takeover schedule, not a scalability or fairness claim for
high-contention deployments. Concurrent-claim rejection is exercised by the test
suite; the fault-injection matrix supplies no contention-regime evidence for P1.

\paragraph{Platform, and a development/measurement split} All measurements were
collected on Ubuntu inside WSL2 on Windows~11, with Docker Desktop port
forwarding in the path, on mains power with sleep and hibernation disabled.
Development and part of the test suite run on the Windows host; \emph{every}
number in \cref{sec:evaluation} comes from the Linux tree against real Redis.
The locally runnable Windows suite and the artifact's Linux suite are checked
separately, while the Linux evidence is the one the measurements rest on.
Absolute latency and throughput are platform-bound; only comparisons
\emph{between} systems sharing this platform are meaningful, which is why
\cref{tab:latency} is presented as increments over a common floor.

\paragraph{The motivating traces come from our own harness} \Cref{sec:motivating}
reproduces failures in baselines we implemented, not in deployed systems. It
establishes that the failure modes are reachable under the stated fault model;
it is not evidence about their frequency in production deployments, and we have
none. A field study of how often such systems meet non-idempotent endpoints
without idempotency keys would strengthen the motivation considerably.

\paragraph{Coverage of the matrix is partial, and the gaps are named} The full
plan is 1\,068 runs; we collected \RunsCollected{}, comprising
\ExecutionsCollected{} executions over \CellsCollected{} cells, chosen to answer
the research questions rather than to fill the grid.
% source: experiments/results/matrix/MANIFEST.md, which lists the run count
% for every cell, and analysis/coverage.json.
The supplementary material enumerates six named gaps --- three since closed and
recorded there with what closing them changed, three still open: the $30\%$
crash-probability regime, the alternative read-back keying, and one run that did
not complete. Two qualifications from that enumeration bear on the main results
and are stated here rather than only there: every collection in this paper comes
from \emph{one host}, so the prevention effect's dependence on host timing is
measured \emph{on} that host rather than established \emph{across} hosts; and
completing a previously partial cell changed its value materially, so a partial
cell should be read as unrepresentative rather than merely imprecise. A rate we
did not measure is absent from the tables rather than interpolated, and the
manifest shipped with the results states the run count for every cell.

\paragraph{One endpoint for the cost measurement} RQ3's cells were collected
against a single endpoint. Overhead should not depend on an endpoint's
reconciliation capability in a crash-free run, since no read-back occurs, but
that is an argument and not a measurement.

\subsection{Conclusion validity}

\paragraph{Every timing number rests on three runs, and one does not survive its
own interval} Each cell behind \cref{tab:deployment} is three crash-free runs of
ten executions, and the four medians there are what the barrier's cost is
computed from. A cluster bootstrap over runs (\cref{sec:eval-deployment}) changes
what can be claimed. Under \texttt{everysec} the barrier's cost is
\BarrierCost{}\,ms with a 95\% interval of
[\BarrierCostLow{}, \BarrierCostHigh{}] --- positive, but pinned only to within
about a factor of four. Under \texttt{always} the point estimate is
\BarrierCostAlways{}\,ms and the interval [\BarrierCostAlwaysLow{},
\BarrierCostAlwaysHigh{}] spans zero: at three runs per arm we cannot show the
barrier costs anything under that policy. An earlier draft quoted a factor
between the two point estimates; it is not supported by these data and has been
removed. The cause is the sample, not the metric: three clusters admit only ten
distinct bootstrap multisets, and the \texttt{always} cell has a heavy upper tail
(p95 \AepAlwaysPninetyfive{}\,ms against a \AepAlwaysMedian{}\,ms median) that
three clusters cannot resolve. Closing this needs more crash-free \emph{runs},
not more executions per run. A deployment that needs a latency budget should
measure its own.

\paragraph{A failure to reject is not a null result here} Two comparisons in
this paper fail to exclude zero for reasons of precision rather than of effect.
The kill-latency attribution above has a half-width wider than the mean it
brackets. So does the capability-class sweep reported in the supplementary
material, where the applied-effect column moved --- contradicting our own
registered prediction --- while the session-clustered interval on the difference
still contains zero because the sessions disagree with one another. In both
cases what we report is a failure to reject at a precision inadequate to the
question, and neither should be read as evidence that the effect is absent.

\paragraph{Timing hygiene, and what it cost us} The host's modern-standby
setting suspended it mid-run during earlier collection, silently inflating
wall-clock durations. We now require both a pre-run declaration that the host
cannot suspend and a per-run wall-versus-monotonic check. The gate is strict
enough that \emph{no} run collected before it existed contributes a timing
number, and it costs us coverage: recovery-latency distributions are not
reported (\cref{sec:eval-rq4}) because the gated crashed runs are too few, and
we prefer the gap to ungated numbers.

\subsection{Where the protocol should not be used}

If the endpoint accepts an idempotency key, or can be enrolled in a
transaction, or the effect can be made idempotent, then a durable-execution
engine gives a strictly stronger guarantee at lower cost and AEP is the wrong
tool. Its cost is real: two fsync barriers, which under \texttt{appendfsync
everysec} dominate the protocol's latency by a factor of roughly
\BarrierToProtocolRatio{} (\cref{tab:latency}) --- a ratio of medians, and one
whose denominator is pinned only to
[\ProtocolMinusBarrierLow{}, \ProtocolMinusBarrierHigh{}]\,ms at three runs per
arm, so the factor is an estimate of an order of magnitude rather than of a
figure. AEP is for the case where the precondition those systems
require cannot be obtained, and its value is a declared residual rather than a
better guarantee.
